\chapter{1993年全国卷（理）}
\section{单选题}

\begin{problem}
如果双曲线的实半轴长为 \(2\)，焦距为 \(6\)，那么该双曲线的离心率为（\quad）

\choices
    {\(\frac{\sqrt{3}}{2}\)}
    {\(\frac{\sqrt{6}}{2}\)}
    {\(\frac{3}{2}\)}
    {\(2\)}
\end{problem}

\begin{answer}
C
\end{answer}

\begin{solution}
设双曲线的实半轴长为 \(a=2\)，焦距为 \(2c=6\)，则 \(c=3\). 因此离心率
\(e=\frac{c}{a}=\frac{3}{2}\)，故选 C.
\end{solution}

\begin{problem}
函数 \(y = \frac{1 - \tan^2 2x}{1 + \tan^2 2x} \) 的最小正周期是（\quad）

\choices
    {\(\frac{\pi}{4}\)}
    {\(\frac{\pi}{2}\)}
    {\(\pi\)}
    {\(2\pi\)}
\end{problem}

\begin{answer}
B
\end{answer}

\begin{solution}
当 \(\tan 2x\) 有意义时，
\(\frac{1-\tan^2 2x}{1+\tan^2 2x}=\cos 4x\). 原式的定义域排除了
\(x=\frac{\pi}{4}+k\frac{\pi}{2}\)（\(k\in\Z\)），这个排除点集的最小正周期也是
\(\frac{\pi}{2}\)，而 \(\cos 4x\) 的最小正周期为 \(\frac{\pi}{2}\)，故选 B.
\end{solution}

\begin{problem}
当圆锥的侧面积和底面积的比值是 \(\sqrt{2}\) 时，圆锥的轴截面顶角是（\quad）

\choices
    {\(45^{\circ}\)}
    {\(60^{\circ}\)}
    {\(90^{\circ}\)}
    {\(120^{\circ}\)}
\end{problem}

\begin{answer}
C
\end{answer}

\begin{solution}
设圆锥的底面半径为 \(r\)，母线长为 \(l\)，轴截面顶角的一半为
\(\theta\). 由侧面积与底面积之比得
\(\frac{\pi rl}{\pi r^2}=\frac{l}{r}=\sqrt{2}\)，所以
\(\sin\theta=\frac{r}{l}=\frac{1}{\sqrt{2}}\). 因为
\(0<\theta<\frac{\pi}{2}\)，故 \(\theta=45^\circ\)，轴截面顶角为
\(2\theta=90^\circ\)，故选 C.
\end{solution}

\begin{problem}
当 \(z = -\frac{1 - \i}{\sqrt{2}}\) 时，\(z^{100} + z^{50} + 1\) 的值等于（\quad）

\choices
    {\(1\)}
    {\(-1\)}
    {\(\i\)}
    {\(-\i\)}
\end{problem}

\begin{answer}
D
\end{answer}

\begin{solution}
\(z=\frac{-1+\i}{\sqrt{2}}=\cos\frac{3\pi}{4}+\i\sin\frac{3\pi}{4}\)，因此
\(z^{50}=\cos\frac{75\pi}{2}+\i\sin\frac{75\pi}{2}=-\i\)，从而
\(z^{100}=-1\). 所以
\(z^{100}+z^{50}+1=-1-\i+1=-\i\)，故选 D.
\end{solution}

\begin{problem}
直线 \(bx + ay = ab\)（\(a < 0\)，\(b < 0\)）的倾斜角是（\quad）

\choices
    {\(\arctan \left(-\frac{b}{a}\right)\)}
    {\(\arctan \left(-\frac{a}{b}\right)\)}
    {\(\pi - \arctan \frac{b}{a}\)}
    {\(\pi - \arctan \frac{a}{b}\)}
\end{problem}

\begin{answer}
C
\end{answer}

\begin{solution}
由直线方程得
\[
y=-\frac{b}{a}x+b.
\]
由于 \(a<0\)，\(b<0\)，所以 \(\frac{b}{a}>0\)，直线的斜率 \(-\frac{b}{a}<0\). 设直线的倾斜角为 \(\alpha\)，则 \(0\leqslant\alpha<\pi\)，且
\[
\tan\alpha=-\frac{b}{a}<0.
\]
因此 \(\frac{\pi}{2}<\alpha<\pi\). 又因为 \(\arctan\frac{b}{a}\in\left(0,\frac{\pi}{2}\right)\)，所以
\[
\alpha=\pi-\arctan\frac{b}{a}.
\]
故选 C.
\end{solution}

\begin{problem}
在直角三角形中两锐角为 \(A\) 和 \(B\)，则 \(\sin A \sin B\)（\quad）

\choices
    {有最大值 \(\frac{1}{2}\) 和最小值 \(0\)}
    {有最大值 \(\frac{1}{2}\)，但无最小值}
    {既无最大值，也无最小值}
    {有最大值 \(1\)，但无最小值}
\end{problem}

\begin{answer}
B
\end{answer}

\begin{solution}
直角三角形的两锐角满足 \(A+B=\frac{\pi}{2}\)，故
\(\sin A\sin B=\sin A\cos A=\frac{1}{2}\sin 2A\). 由于
\(0<2A<\pi\)，有 \(0<\sin A\sin B\leqslant\frac{1}{2}\). 当
\(A=\frac{\pi}{4}\) 时取得最大值 \(\frac{1}{2}\)，而 \(A\) 不能取 \(0\) 或
\(\frac{\pi}{2}\)，所以最小值 \(0\) 不能取得，故选 B.
\end{solution}

\begin{problem}
在各项均为正数的等比数列 \(\{a_{n}\}\) 中，若 \(a_{5}a_{6} = 9\)，则 \(\log_{3} a_{1} + \log_{3} a_{2} + \dots + \log_{3} a_{10}\) 的值为（\quad）

\choices
    {\(12\)}
    {\(10\)}
    {\(8\)}
    {\(2 + \log_{3} 5\)}
\end{problem}

\begin{answer}
B
\end{answer}

\begin{solution}
设等比数列的公比为 \(q\). 对 \(i=1,2,\dots,5\)，数列通项的指数相加为
\((i-1)+(10-i)=9\)，所以
\(a_i a_{11-i}=a_5a_6=9\). 因而
\[
\sum_{i=1}^{10}\log_3 a_i
=\sum_{i=1}^{5}\log_3(a_i a_{11-i})
=5\log_3 9=10.
\]
故选 B.
\end{solution}

\begin{problem}
\(F(x) = \left(1 + \frac{2}{2^{x} - 1}\right)f(x)\)（\(x \neq 0\)）是偶函数，且 \(f(x)\) 不恒等于零，则 \(f(x)\)（\quad）

\choices
    {是奇函数}
    {是偶函数}
    {可能是奇函数也可能是偶函数}
    {不是奇函数也不是偶函数}
\end{problem}

\begin{answer}
A
\end{answer}

\begin{solution}
令
\(g(x)=1+\frac{2}{2^x-1}=\frac{2^x+1}{2^x-1}\). 当 \(x\neq0\) 时，
\(g(-x)=-g(x)\)，且 \(g(x)\neq0\). 由 \(F(x)\) 为偶函数，得
\(-g(x)f(-x)=g(x)f(x)\)，约去 \(g(x)\) 后得到
\(f(-x)=-f(x)\). 因此 \(f(x)\) 是奇函数，故选 A.
\end{solution}

\begin{problem}
曲线的参数方程为 \(\begin{cases}
x = 3t^2 + 2,\\
y = t^2 - 1,
\end{cases}\) （\(0 \leqslant t \leqslant 5\)），则曲线是（\quad）

\choices
    {线段}
    {双曲线的一支}
    {圆弧}
    {射线}
\end{problem}

\begin{answer}
A
\end{answer}

\begin{solution}
消去参数得 \(x=3(t^2-1)+5=3y+5\). 又因为
\(0\leqslant t\leqslant5\)，所以 \(0\leqslant t^2\leqslant25\)，从而
\(-1\leqslant y\leqslant24\). 因此曲线是直线
\(x=3y+5\) 上从 \((2,-1)\) 到 \((77,24)\) 的线段，故选 A.
\end{solution}

\begin{problem}
若 \(a\)，\(b\) 是任意实数，且 \(a > b\)，则（\quad）

\choices
    {\(a^2 > b^2\)}
    {\(\frac{b}{a} < 1\)}
    {\(\lg (a - b) > 0\)}
    {\(\left(\frac{1}{2}\right)^a < \left(\frac{1}{2}\right)^b\)}
\end{problem}

\begin{answer}
D
\end{answer}

\begin{solution}
由 \(a>b\) 只能得到 \(a-b>0\). 例如取 \((a,b)=(1,-2)\)，则
\(a^2>b^2\) 不成立；取 \((a,b)=(-1,-2)\)，则
\(\frac{b}{a}<1\) 不成立；取 \((a,b)=(1,\frac{1}{2})\)，则
\(\lg(a-b)>0\) 不成立. 由于底数
\(\frac{1}{2}\) 在 \((0,1)\) 内，指数函数严格递减，所以
\(a>b\) 时有
\(\left(\frac{1}{2}\right)^a<\left(\frac{1}{2}\right)^b\)，故选 D.
\end{solution}

\begin{problem}
已知集合 \(E = \{\theta | \cos \theta < \sin \theta, 0 \leqslant \theta \leqslant 2\pi\}\)，\(F = \{\theta | \tan \theta < \sin \theta\}\)，那么 \(E \cap F\) 为区间（\quad）

\choices
    {\(\left(\frac{\pi}{2},\pi\right)\)}
    {\(\left(\frac{\pi}{4},\frac{3\pi}{4}\right)\)}
    {\(\left(\pi,\frac{3\pi}{2}\right)\)}
    {\(\left(\frac{3\pi}{4},\frac{5\pi}{4}\right)\)}
\end{problem}

\begin{answer}
A
\end{answer}

\begin{solution}
由 \(\sin\theta-\cos\theta>0\) 及 \(0\leqslant\theta\leqslant2\pi\)，得
\(E=\left(\frac{\pi}{4},\frac{5\pi}{4}\right)\). 在正切有意义时，
\[
\tan\theta-\sin\theta=\frac{\sin\theta(1-\cos\theta)}{\cos\theta}.
\]
在 \(\left(\frac{\pi}{4},\frac{\pi}{2}\right)\) 内该式为正，在
\(\left(\frac{\pi}{2},\pi\right)\) 内为负，在
\(\left(\pi,\frac{5\pi}{4}\right)\) 内为正. 因此
\(E\cap F=\left(\frac{\pi}{2},\pi\right)\)，故选 A.
\end{solution}

\begin{problem}
一动圆与两圆：\(x^{2} + y^{2} = 1\) 和 \(x^{2} + y^{2} - 8x + 12 = 0\) 都外切，则动圆圆心的轨迹为（\quad）

\choices
    {抛物线}
    {圆}
    {双曲线的一支}
    {椭圆}
\end{problem}

\begin{answer}
C
\end{answer}

\begin{solution}
设动圆圆心为 \(M\)，半径为 \(r\)，两个定圆的圆心分别为
\(O_1=(0,0)\)、\(O_2=(4,0)\)，半径分别为 \(1\)、\(2\). 两定圆圆心距为
\(O_1O_2=4\). 由外切条件，
\(MO_1=r+1\)，\(MO_2=r+2\)，所以
\(MO_2-MO_1=1\). 动点到两个定点的距离之差为定值，故其轨迹为双曲线的一支，故选 C.
\end{solution}

\begin{problem}
若正棱锥的底面边长与侧棱长相等，则该棱锥一定不是（\quad）

\choices
    {三棱锥}
    {四棱锥}
    {五棱锥}
    {六棱锥}
\end{problem}

\begin{answer}
D
\end{answer}

\begin{solution}
设正 \(n\) 边形底面边长为 \(s\)，其外接圆半径为
\(R=\frac{s}{2\sin(\pi/n)}\). 侧棱长等于 \(s\) 时，棱锥高满足
\(h^2=s^2-R^2\)，因此必须有 \(R<s\). 当 \(n=6\) 时
\(R=s\)，从而 \(h=0\)，不能构成棱锥；而 \(n=3,4,5\) 时均有 \(R<s\). 故一定不是六棱锥，选 D.
\end{solution}

\begin{problem}
如果圆柱轴截面的周长 \(l\) 为定值，那么圆柱体积的最大值是（\quad）

\choices
    {\(\left(\frac{l}{6}\right)^3\pi\)}
    {\(\left(\frac{l}{3}\right)^3\pi\)}
    {\(\left(\frac{l}{4}\right)^3\pi\)}
    {\(\frac{1}{4}\left(\frac{l}{4}\right)^3\pi\)}
\end{problem}

\begin{answer}
A
\end{answer}

\begin{solution}
设圆柱底面半径为 \(r\)，高为 \(h\). 轴截面周长为 \(l\)，故
\(4r+2h=l\)，即 \(h=\frac{l}{2}-2r\). 于是
\[
V=\pi r^2h=\pi\left(\frac{l}{2}r^2-2r^3\right),
\qquad
V'=\pi r(l-6r).
\]
在 \(0<r<\frac{l}{4}\) 内，唯一临界点为 \(r=\frac{l}{6}\)，此时
\(h=\frac{l}{6}\)，且两端体积趋于零，所以体积最大值为
\(\pi\left(\frac{l}{6}\right)^3\)，故选 A.
\end{solution}

\begin{problem}
由 \((\sqrt{3}x + \sqrt[3]{2})^{100}\) 展开所得的 \(x\) 的多项式中，系数为有理数的共有（\quad）

\choices
    {\(50\) 项}
    {\(17\) 项}
    {\(16\) 项}
    {\(15\) 项}
\end{problem}

\begin{answer}
B
\end{answer}

\begin{solution}
展开式中 \(x^k\) 的系数为
\(\mathrm C_{100}^{k}3^{k/2}2^{(100-k)/3}\)，其中 \(0\leqslant k\leqslant100\). 该系数为有理数的条件是
\(k\) 为偶数且 \(100-k\) 是 \(3\) 的倍数，即 \(k\equiv4\pmod 6\). 所有可能的 \(k\) 为
\(4,10,16,\dots,100\)，共有
\(\frac{100-4}{6}+1=17\) 个，故选 B.
\end{solution}

\begin{problem}
设 \(a\)，\(b\)，\(c\) 都是正数，且 \(3a = 4b = 6c\)，那么（\quad）

\choices
    {\(\frac{1}{c} = \frac{1}{a} + \frac{1}{b}\)}
    {\(\frac{2}{c} = \frac{2}{a} +\frac{1}{b}\)}
    {\(\frac{1}{c} = \frac{2}{a} +\frac{2}{b}\)}
    {\(\frac{2}{c} = \frac{1}{a} +\frac{2}{b}\)}
\end{problem}

\begin{answer}
无正确选项
\end{answer}

\begin{solution}
令 \(3a=4b=6c=k\)，其中 \(k>0\)，则
\[
a=\frac{k}{3},\qquad b=\frac{k}{4},\qquad c=\frac{k}{6},
\]
从而
\[
\frac{1}{a}=\frac{3}{k},\qquad \frac{1}{b}=\frac{4}{k},\qquad \frac{1}{c}=\frac{6}{k}.
\]
逐项检验可得：选项 A 的左右两边分别为 \(\frac{6}{k}\) 和 \(\frac{7}{k}\)；选项 B 的左右两边分别为 \(\frac{12}{k}\) 和 \(\frac{10}{k}\)；选项 C 的左右两边分别为 \(\frac{6}{k}\) 和 \(\frac{14}{k}\)；选项 D 的左右两边分别为 \(\frac{12}{k}\) 和 \(\frac{11}{k}\).

因此，按原题面，四个选项均不成立，题目没有可选的正确答案.
\end{solution}

\begin{problem}
同室四人各写一张贺年卡，先集中起来，然后每人从中拿一张别人送出的贺年卡，则四张贺年卡不同的分配方式有（\quad）

\choices
    {\(6\) 种}
    {\(9\) 种}
    {\(11\) 种}
    {\(23\) 种}
\end{problem}

\begin{answer}
B
\end{answer}

\begin{solution}
同室四人的分配要求每个人都不能拿到自己写的贺年卡，因此这是四个元素的错排问题. 由容斥原理，错排数为
\[
4!-\mathrm C_4^1 3!+\mathrm C_4^2 2!-\mathrm C_4^3 1!+\mathrm C_4^4 0!=24-24+12-4+1=9.
\]
故选 B.
\end{solution}

\begin{problem}
已知异面直线 \(a\) 与 \(b\) 所成的角为 \(50^{\circ}\)，\(P\) 为空间上一定点，则过点 \(P\) 且与 \(a\)，\(b\) 所成的角都是 \(30^{\circ}\) 的直线有且仅有（\quad）

\choices
    {\(1\) 条}
    {\(2\) 条}
    {\(3\) 条}
    {\(4\) 条}
\end{problem}

\begin{answer}
B
\end{answer}

\begin{solution}
取两条异面直线的单位方向向量为
\(\bs u=(\cos25^\circ,\sin25^\circ,0)\) 和
\(\bs v=(\cos25^\circ,-\sin25^\circ,0)\). 设所求直线的单位方向向量为
\(\bs w=(x,y,z)\). 因为直线所成的角为 \(30^\circ\)，有
\(\lvert\bs w\cdot\bs u\rvert=\lvert\bs w\cdot\bs v\rvert=\cos30^\circ\).

若两个内积同号，取同为 \(\cos30^\circ\)（取负号只会把方向向量整体反向），两式相减得 \(y=0\)，再由其中一式得
\(x=\frac{\cos30^\circ}{\cos25^\circ}<1\). 因而
\(z=\pm\sqrt{1-x^2}\)，得到两条不同的直线.

若两个内积异号，两式相加或相减可得
\(x=0\)，且 \(\lvert y\rvert=\frac{\cos30^\circ}{\sin25^\circ}>1\)，这与
\(x^2+y^2+z^2=1\) 矛盾. 所以所求直线有且仅有 \(2\) 条，故选 B.
\end{solution}

\section{填空题}

\begin{problem}
抛物线 \(y^{2} = 4x\) 的弦 \(AB\) 垂直于 \(x\) 轴，若 \(AB\) 的长为 \(4\sqrt{3}\)，则焦点到 \(AB\) 的距离为\fillinblank{}.
\end{problem}

\begin{answer}
2
\end{answer}

\begin{solution}
设弦 \(AB\) 所在的直线为 \(x=x_0\). 代入抛物线
\(y^2=4x\)，弦端点的纵坐标为 \(\pm2\sqrt{x_0}\)，所以
\(AB=4\sqrt{x_0}=4\sqrt3\)，得 \(x_0=3\). 抛物线
\(y^2=4x\) 的焦点为 \((1,0)\)，故焦点到直线 \(x=3\) 的距离为 \(2\).
\end{solution}

\begin{problem}
在半径为 \(30 \mathrm{~m}\) 的圆形广场中央上空，设置一个照明光源，射向地面的光呈圆锥形，且其轴截面顶角为 \(120^{\circ}\). 若要光源恰好照亮整个广场，则其高度应为\fillinblank{} \(\mathrm{m}\)（精确到 \(0.1 \mathrm{~m}\)）.
\end{problem}

\begin{answer}
17.3
\end{answer}

\begin{solution}
设光源离地面的高度为 \(h\). 轴截面顶角为 \(120^\circ\)，所以其半角为
\(60^\circ\). 在轴截面的直角三角形中，
\(\tan60^\circ=\frac{30}{h}\)，从而
\(h=\frac{30}{\sqrt3}=10\sqrt3\approx17.3205\). 按题意精确到
\(0.1\mathrm{~m}\)，得 \(17.3\).
\end{solution}

\begin{problem}
在 \(50\) 件产品中有 \(4\) 件是次品，从中任意抽出 \(5\) 件，至少有 \(3\) 件是次品的抽法共\fillinblank{}\(\text{种}\)（用数字作答）.
\end{problem}

\begin{answer}
4186
\end{answer}

\begin{solution}
至少抽到 \(3\) 件次品包括恰有 \(3\) 件和恰有 \(4\) 件两种情形. 因此抽法数为
\[
\mathrm C_4^3\mathrm C_{46}^2+\mathrm C_4^4\mathrm C_{46}^1
=4\times1035+46=4186.
\]
\end{solution}

\begin{problem}
建造一个容积为 \(8 \mathrm{~m}^{3}\)，深为 \(2 \mathrm{~m}\) 的长方体无盖水池，如果池底和池壁的造价每平方米分别为 \(120\text{ 元}\) 和 \(80\text{ 元}\)，那么水池的最低总造价为\fillinblank{}\(\text{元}\).
\end{problem}

\begin{answer}
1760
\end{answer}

\begin{solution}
设池底长、宽分别为 \(x\mathrm{~m}\)、\(y\mathrm{~m}\). 由容积条件
\(2xy=8\)，得 \(xy=4\). 池底造价为 \(120xy=480\) 元，四面池壁面积为
\(4x+4y\)，池壁造价为 \(320(x+y)\) 元. 因此总造价
\(C=480+320(x+y)\). 由基本不等式
\(x+y\geqslant2\sqrt{xy}=4\)，当且仅当 \(x=y=2\) 时取等号，所以
\(C_{\min}=480+320\times4=1760\) 元.
\end{solution}

\begin{problem}
设 \(f(x) = 4^{x} - 2^{x + 1}\)，则 \(f^{-1}(0) =\)\fillinblank{}.
\end{problem}

\begin{answer}
1
\end{answer}

\begin{solution}
令 \(t=2^x>0\)，则
\(f(x)=2^{2x}-2^{x+1}=t^2-2t=t(t-2)\). 因为 \(t>0\)，所以
\(f(x)=0\) 当且仅当 \(t=2\)，即 \(2^x=2\)，从而 \(x=1\). 因此
\(f^{-1}(0)=1\).
\end{solution}

\begin{problem}
已知等差数列 \(\{a_{n}\}\) 的公差 \(d > 0\)，首项 \(a_{1} > 0\)，\(S_{n} = \sum_{i=1}^{n} \frac{1}{a_{i} a_{i+1}}\)，则 \(\displaystyle\lim_{n \to \infty} S_{n} =\)\fillinblank{}.
\end{problem}

\begin{answer}
\(\frac{1}{a_1d}\)
\end{answer}

\begin{solution}
设 \(a_i=a_1+(i-1)d\)，则 \(a_{i+1}-a_i=d\). 因为 \(a_1>0\)，\(d>0\)，所以每一项都不为零，并且
\[
\frac{1}{a_i a_{i+1}}=\frac{1}{d}\left(\frac{1}{a_i}-\frac{1}{a_{i+1}}\right).
\]
于是
\[
S_n=\frac{1}{d}\sum_{i=1}^{n}\left(\frac{1}{a_i}-\frac{1}{a_{i+1}}\right)=\frac{1}{d}\left(\frac{1}{a_1}-\frac{1}{a_{n+1}}\right).
\]
又因为 \(a_{n+1}=a_1+nd\to+\infty\)，所以 \(\frac{1}{a_{n+1}}\to0\). 因而
\[
\lim_{n\to\infty}S_n=\frac{1}{a_1d}.
\]
故填 \(\frac{1}{a_1d}\).
\end{solution}

\section{解答题}

\begin{problem}
解不等式：\(2 + \log_{\frac{1}{2}}(5 - x) + \log_{2}\frac{1}{x} > 0\).
\end{problem}

\begin{solution}
由对数的真数必须为正，得
\[
5-x>0,\qquad x>0,
\]
所以定义域为 \(0<x<5\). 在此范围内，利用换底关系和倒数的对数性质，原不等式等价于
\[
2-\log_{2}(5-x)-\log_{2}x>0,
\]
即
\[
\log_{2}\bigl(x(5-x)\bigr)<2.
\]
由于底数 \(2>1\)，这又等价于
\[
x(5-x)<4,
\]
从而
\[
x^2-5x+4>0,
\qquad
(x-1)(x-4)>0.
\]
结合定义域 \(0<x<5\)，得
\[
x\in(0,1)\cup(4,5).
\]
故原不等式的解集为 \((0,1)\cup(4,5)\).
\end{solution}

\begin{problem}
如图，\(A_{1}B_{1}C_{1}-ABC\) 是直三棱柱，过点 \(A_{1}\)，\(B\)，\(C_{1}\) 的平面和平面 \(ABC\) 的交线记作 \(l\).

\begin{enumerate}
\item 判定直线 \(A_{1}C_{1}\) 和 \(l\) 的位置关系，并加以证明；
\item 若 \(A_{1}A = 1\)，\(AB = 4\)，\(BC = 3\)，\(\angle ABC = 90^{\circ}\)，求顶点 \(A_{1}\) 到直线 \(l\) 的距离.
\end{enumerate}

\bitmapfigure[max width=0.42\linewidth]{img/1993/national_science/q26_fig1.png}
\end{problem}

\begin{solution}
\begin{enumerate}
\item 建立空间直角坐标系，使 \(B=(0,0,0)\)，\(BA\) 为 \(x\) 轴正方向，\(BC\) 为 \(y\) 轴正方向，且 \(AA_1\) 为 \(z\) 轴方向. 由题设可取
\[
A=(4,0,0),\qquad C=(0,3,0),\qquad A_1=(4,0,1),\qquad C_1=(0,3,1).
\]
于是直线 \(A_1C_1\) 的方向向量为 \((-4,3,0)\). 过 \(A_1\)，\(B\)，\(C_1\) 的平面由向量 \((4,0,1)\) 和 \((0,3,1)\) 张成，其方程为
\[
3x+4y-12z=0.
\]
它与平面 \(ABC\)（即 \(z=0\)）的交线 \(l\) 满足
\[
3x+4y=0.
\]
因此 \(l\) 的一个方向向量为 \((4,-3,0)\)，与 \((-4,3,0)\) 平行. 所以 \(A_1C_1\parallel l\).

\item 直线 \(l\) 经过 \(B\)，方向向量可取 \(\bs u=(4,-3,0)\). 点 \(A_1\) 的位置向量为 \(\overrightarrow{BA_1}=(4,0,1)\). 点到直线的距离为
\[
\frac{\left\|\overrightarrow{BA_1}\times\bs u\right\|}{\left\|\bs u\right\|}
=\frac{\left\|(3,4,-12)\right\|}{5}
=\frac{13}{5}.
\]
故顶点 \(A_1\) 到直线 \(l\) 的距离为 \(\frac{13}{5}\).
\end{enumerate}
\end{solution}

\begin{problem}
在面积为 \(1\) 的 \(\triangle PMN\) 中，\(\tan M = \frac{1}{2}\)，\(\tan N = -2\)，建立适当的坐标系，求出以 \(M\)，\(N\) 为焦点且过点 \(P\) 的椭圆方程.
\end{problem}

\begin{solution}
建立直角坐标系，使 \(M=(-c,0)\)，\(N=(c,0)\)，并取点 \(P\) 在 \(x\) 轴上方. 记三角形的内角仍为 \(M\)，\(N\)，\(P\). 因为 \(M\) 为锐角，\(N\) 为钝角，且
\[
\tan M=\frac{1}{2},\qquad \tan N=-2=-\frac{1}{\tan M},
\]
所以 \(N=\frac{\pi}{2}+M\)，从而
\[
P=\pi-M-N=\frac{\pi}{2}-2M.
\]
由 \(\tan M=\frac{1}{2}\)，得 \(\sin M=\frac{1}{\sqrt{5}}\)，\(\cos M=\frac{2}{\sqrt{5}}\)，\(\sin N=\frac{2}{\sqrt{5}}\). 又有 \(\sin P=\cos 2M=\frac{3}{5}\)，\(\cos P=\sin 2M=\frac{4}{5}\).

由正弦定理，\(PM:PN=\sin N:\sin M=2:1\). 设 \(PN=t\)，则 \(PM=2t\). 由三角形面积为 \(1\)，有
\[
1=\frac{1}{2}\cdot PM\cdot PN\cdot\sin P=t^2\cdot\frac{3}{5},
\]
所以 \(t^2=\frac{5}{3}\). 再由余弦定理，
\[
MN^2=PM^2+PN^2-2\cdot PM\cdot PN\cdot\cos P
=5t^2-\frac{16}{5}t^2=3,
\]
即 \(MN=\sqrt{3}\). 另一方面，点 \(P\) 在所求椭圆上，所以椭圆的长轴长满足
\[
2a=PM+PN=3t=\sqrt{15}.
\]
因此 \(a=\frac{\sqrt{15}}{2}\)，而两个焦点间距为 \(2c=MN=\sqrt{3}\)，故 \(c=\frac{\sqrt{3}}{2}\). 于是
\[
b^2=a^2-c^2=\frac{15}{4}-\frac{3}{4}=3.
\]
所求椭圆方程为
\[
\frac{x^2}{15/4}+\frac{y^2}{3}=1.
\]
\end{solution}

\begin{problem}
设复数 \(z = \cos \theta + \i\sin \theta\)（\(0 < \theta < \pi\)），\(\omega = \frac{1 - \left( \overline{z} \right)^4}{1 + z^4}\)，已知 \(|\omega| = \frac{\sqrt{3}}{3}\)，\(\arg \omega = \frac{\pi}{2}\)，求 \(\theta\).
\end{problem}

\begin{solution}
由 \(z=\e^{\i\theta}\)，有
\[
1-\left(\overline z\right)^4=1-\e^{-4\i\theta}=2\i\e^{-2\i\theta}\sin 2\theta,
\qquad
1+z^4=1+\e^{4\i\theta}=2\e^{2\i\theta}\cos 2\theta.
\]
题设要求 \(\omega\) 有意义，所以 \(\cos 2\theta\ne0\). 因而
\[
\omega=\i\e^{-4\i\theta}\tan 2\theta
=\tan 2\theta\sin 4\theta+\i\tan 2\theta\cos 4\theta.
\]
由 \(|\omega|=\frac{\sqrt{3}}{3}>0\)，知 \(\tan 2\theta\ne0\). 又 \(\arg\omega=\frac{\pi}{2}\) 意味着 \(\omega\) 为正纯虚数，所以其实部为零，从而
\[
\tan 2\theta\sin 4\theta=0,
\qquad
\sin 4\theta=0.
\]
于是 \(\theta=\frac{k\pi}{4}\). 结合 \(0<\theta<\pi\)，只有 \(k=1\)，\(2\)，\(3\) 需要讨论. 当 \(k=1\)，\(3\) 时，\(\cos 2\theta=0\)，原式分母为零；当 \(k=2\) 时，\(\tan 2\theta=0\)，从而 \(\omega=0\)，与 \(|\omega|=\frac{\sqrt{3}}{3}\) 矛盾.

所以按原题面条件，不存在满足条件的 \(\theta\).
\end{solution}

\begin{problem}
已知关于 \(x\) 的实系数二次方程 \(x^{2} + ax + b = 0\) 有两个实数根 \(\alpha\)，\(\beta\). 证明：

\begin{enumerate}
\item 如果 \(|\alpha| < 2\)，\(|\beta| < 2\)，那么 \(2|\alpha| < 4 + b\) 且 \(|b| < 4\)；
\item 如果 \(2|\alpha| < 4 + b\) 且 \(|b| < 4\)，那么 \(|\alpha| < 2\)，\(|\beta| < 2\).
\end{enumerate}
\end{problem}

\begin{solution}
按原题面，题中两条命题均不成立，因此不能给出有效的证明. 由韦达定理，\(b=\alpha\beta\)，可以直接构造反例.

\begin{enumerate}
\item 取 \(\alpha=\frac{3}{2}\)，\(\beta=-\frac{19}{10}\)，则 \(|\alpha|<2\)，\(|\beta|<2\)，且对应的实系数二次方程为
\[
x^2-\left(\alpha+\beta\right)x+\alpha\beta=x^2+\frac{2}{5}x-\frac{57}{20}=0.
\]
但是
\[
2|\alpha|=3>\frac{23}{20}=4+b,
\qquad
|b|=\frac{57}{20}<4.
\]
所以第（1）条中的第一个不等式不成立.

\item 取 \(\alpha=\frac{5}{2}\)，\(\beta=\frac{3}{2}\)，则对应的实系数二次方程为
\[
x^2-4x+\frac{15}{4}=0.
\]
此时
\[
2|\alpha|=5<\frac{31}{4}=4+b,
\qquad
|b|=\frac{15}{4}<4,
\]
但 \(|\alpha|=\frac{5}{2}>2\). 所以第（2）条的结论也不成立.
\end{enumerate}

因此，第29题的原题面存在客观矛盾；本稿保留原题面，不能将其改写成另一道题.
\end{solution}
